% 大学物理核心公式与考点汇总 - XeLaTeX无错优化版（支持五级标题TOC）
% 层级：chapter(1) > section(2) > subsection(3) > subsubsection(4) > subsubsubsection(5)
% 编号示例：2.1.2.5.1，适配MiKTeX/TeXLive Windows
\documentclass[12pt,a4paper,oneside]{report}

% ======================== 核心宏包 ========================
\usepackage{geometry}
\usepackage{microtype}  % 优化字符间距，减少Underfull/Overfull警告
\usepackage{needspace}

% 中文与字体【修复顺序冲突】
\usepackage[UTF8, fontset=windows]{ctex}
\usepackage{fontspec}
\setmainfont{Times New Roman}
\setsansfont{Microsoft YaHei}

% 数学与格式
\usepackage{amsmath,amssymb,amsfonts}
\usepackage{upgreek}
\usepackage{xcolor}
\usepackage{colortbl}
\usepackage{booktabs}
\usepackage{multirow}
\usepackage{makecell}
\usepackage{siunitx}
\usepackage{titlesec}
\usepackage{titletoc}
\usepackage{tcolorbox}
\usepackage{fancyhdr}
\usepackage{enumitem}

% 超链接顺序规范【修复hypcap必须在hyperref之后加载】
\usepackage[
    bookmarks=true,unicode=true,colorlinks=true,
    linkcolor=mainblue,citecolor=mainblue,urlcolor=mainblue,
    pdfborder={0 0 0}
]{hyperref}
\usepackage{hypcap}

% 告诉 hyperref 五级标题的目录层级，消除 bookmark 警告
\makeatletter
\providecommand\toclevel@subsubsubsection{5}
\makeatother

% ======================== 基础页面设置 ========================
\geometry{a4paper,left=2.5cm,right=2.5cm,top=2.5cm,bottom=2.5cm,headheight=15pt}

% 全局颜色定义
\definecolor{mainblue}{RGB}{41, 90, 141}
\definecolor{secondcyan}{RGB}{26, 117, 117}
\definecolor{thirdteal}{RGB}{27, 120, 101}
\definecolor{highlight}{RGB}{255, 240, 200}
\definecolor{tablehead}{RGB}{230, 240, 250}
\definecolor{textgray}{RGB}{60, 60, 60}

% 页眉页脚全局样式
\pagestyle{fancy}
\fancyhf{}
\fancyhead[L]{\color{mainblue}\bfseries\small 大学物理核心公式与考点汇总}
\fancyhead[R]{\color{mainblue}\thepage}
\renewcommand{\headrulewidth}{0.5pt}
\renewcommand{\headrule}{\color{mainblue}\hrule width\headwidth height\headrulewidth}
% 第一章首页plain页眉统一
\fancypagestyle{plain}{
    \fancyhf{}
    \fancyhead[L]{\color{mainblue}\bfseries\small 大学物理核心公式与考点汇总}
    \fancyhead[R]{\color{mainblue}\thepage}
    \renewcommand{\headrulewidth}{0.5pt}
}

% 各级标题垂直间距紧凑优化（新增五级标题间距）
\titlespacing*{\section}{0pt}{6pt plus 2pt minus 2pt}{3pt plus 2pt minus 2pt}
\titlespacing*{\subsection}{0pt}{4pt plus 2pt minus 2pt}{2pt plus 2pt minus 2pt}
\titlespacing*{\subsubsection}{0pt}{3pt plus 1pt minus 1pt}{1pt plus 1pt minus 1pt}
\titlespacing*{\subsubsubsection}{0pt}{2pt plus 1pt minus 1pt}{1pt plus 1pt minus 1pt}

% 枚举列表全局紧凑配置
\setlist{
    itemsep=0.1ex,
    parsep=0.2ex,
    leftmargin=1.5em,
    labelsep=0.4em,
    topsep=0.2ex,
    partopsep=0.1ex,
    font=\small
}
\setlist[itemize,2]{leftmargin=1.8em,label=$\cdot$,itemsep=0.05ex,parsep=0.1ex}

% 数学公式上下间距压缩
\numberwithin{equation}{section}
\allowdisplaybreaks
\setlength{\abovedisplayskip}{3pt}
\setlength{\belowdisplayskip}{3pt}
\setlength{\abovedisplayshortskip}{2pt}
\setlength{\belowdisplayshortskip}{2pt}
\makeatletter
\renewcommand{\@mathmargin}{0pt}
\makeatother

% 整体排版容错、缩进、段落间距
\setlength{\emergencystretch}{2em}
\sloppy
\setlength{\parindent}{0pt}
\setlength{\parskip}{0.8ex}
\setlength{\tabcolsep}{4pt}

% ======================== 自定义快捷命令 ========================
% 高亮核心公式框
\newcommand{\keyformula}[1]{%
    \needspace{2\baselineskip}%
    \begin{tcolorbox}[
        colback=highlight,colframe=white,boxrule=0pt,arc=0pt,
        left=6pt,right=6pt,top=4pt,bottom=4pt,center upper,
        boxsep=0pt
    ]
    $\displaystyle #1$
    \end{tcolorbox}%
    \vspace{0.2em}
}

% 符号说明命令
\newcommand{\symb}[2]{\textbf{#1}：\parbox[t]{10cm}{\raggedright #2}\\[0.3ex]}

% 核心考点蓝色边框提示框
\newcommand{\corepoint}[1]{%
    \begin{tcolorbox}[
        colback=white,colframe=mainblue,boxrule=0.5pt,arc=2pt,
        left=6pt,right=6pt,top=2pt,bottom=2pt,center upper,
        boxsep=0pt
    ]
    \textbf{\color{mainblue}核心考点：#1}
    \end{tcolorbox}%
    \vspace{0.2em}
}

% ======================== 【核心：完整定义五级标题 subsubsubsection + TOC目录渲染】 ========================
\makeatletter
% 1. 创建五级计数器，依附于四级 subsubsection
\newcounter{subsubsubsection}[subsubsection]
% 2. 五级编号格式：2.1.2.5.1 链式编号
\renewcommand{\thesubsubsubsection}{\thesubsubsection.\arabic{subsubsubsection}}

% 3. 定义五级标题命令 \subsubsubsection{标题名称}，层级编号=5
\newcommand{\subsubsubsection}{\@startsection
    {subsubsubsection}
    {5}   % 层级编号改为5，对应五级标题
    {\z@}
    {3.25ex \@plus1ex \@minus.2ex}
    {1.5ex \@plus.2ex}
    {\normalfont\normalsize\bfseries}
}

% 4. 重定义目录内部命令 \l@subsubsubsection，让tocdepth=5时正常生成五级目录条目
\def\l@subsubsubsection#1#2{%
  \ifnum \c@tocdepth >4  % tocdepth>4 才渲染五级
    \addpenalty{-\@highpenalty}%
    \vskip 0.4em \@plus\p@
    \setlength\@tempdima{16em}% 五级目录缩进宽度，按需修改
    \begingroup
      \parindent \z@ \rightskip \@pnumwidth
      \parfillskip -\@pnumwidth
      \leavevmode
      \advance\leftskip\@tempdima
      \hskip -\leftskip
      #1\nobreak\hfil \nobreak\hb@xt@\@pnumwidth{\hss #2}\par
      \penalty\@highpenalty
    \endgroup
  \fi
}
\makeatother

% ======================== 目录深度关键配置（开启五级显示） ========================
\setcounter{secnumdepth}{5}   % 正文允许编号到第5级标题
\setcounter{tocdepth}{5}      % 目录TOC渲染显示到第5级标题

% ======================== 文档入口 ========================
\begin{document}

% 封面页
% 封面定义
\newcommand{\makeCover}
{
  \begin{titlepage}
    % 临时调整封面边距
    \newgeometry{left=2cm, right=2cm, top=2cm, bottom=2cm}
    
    \centering
    \vspace*{1cm}
    
    % 牛顿墓志铭（诗意改写版）
    {\Large \itshape 他以近乎神性的心智，\\以数学为炬，\\照亮行星运转、彗星行迹与海洋潮汐。\\ 艾萨克·牛顿的墓志铭}\\[0.5cm]
    \rule{\textwidth}{1pt}\\[0.5cm]
    
    % 主标题
    {\Huge \bfseries \color{mainblue} 大学物理}\\[8pt]
    {\Huge \bfseries \color{mainblue} 核心公式与考点汇总}\\[12pt]
    {\Large \itshape —— 简洁版复习讲义 ——}\\[1cm]
    
    % 分隔线
    \rule{0.6\textwidth}{0.5pt}\\[1cm]
    
    % ========== 修改点：更新内容概要 ==========
    {\Large \textbf{内容概要：} 力学 | 电磁学 | 振动\&波动}\\[8pt]
    {\Large \textbf{适用场景：} 课程复习 · 公式速查}\\[1.5cm]
    
    % 整理者信息
    {\Large \textbf{整理者：}浮萍}\\[6pt]
    {\large \textbf{联系邮箱：}\href{mailto:3071974740@qq.com}{3071974740@qq.com}}\\[6pt]
    {\large \textbf{GitHub：}\url{https://github.com/Duckweed-yhb}}\\[6pt]
    {\Large 整理日期：\today}\\[1cm]
    
    \rule{0.6\textwidth}{0.5pt}\\[1cm]
    
    % 底部信息
    {\large \textbf{版本：} v1.0 简洁版}\\[6pt]
    {\small \itshape Copyright © \the\year\ 浮萍 · 禁止商用 · 欢迎交流}\\[0.5cm]
    
    \rule{\textwidth}{1pt}
    \vfill
    
    % 恢复原始页面布局
    \restoregeometry
  \end{titlepage}
}

% 生成封面
\makeCover
\clearpage
\thispagestyle{fancy}

% 自动生成目录（两次XeLaTeX编译生效五级目录）
\tableofcontents
\clearpage

% ========== 第一篇 力学 ==========
\part{第一篇 力学}
\chapter{第一章 质点运动学}

\section{运动时空测量}
    \subsection{物质与运动}
    \begin{enumerate}
        \item 质点运动学：描述机械运动物体空间位置随时间的变化关系，不探究运动产生与改变的原因。
        \item 物质定义：客观存在，能量的表达；以\textbf{实物}和\textbf{场}两种基本形式存在，运动是物质固有属性，场是更基本的物质形态。
    \end{enumerate}

    \subsection{测量的标准}
    \begin{enumerate}
        \item 物理量测量规则：将物理量与公认标准比对，单位是物理量的命名，标准为精确不变的基准。
        \item 标准选取原则：易于获取、恒定不变。
    \end{enumerate}

    \subsection{国际单位制与量纲}
    \subsubsection{SI国际单位制7个基本量}
    \begin{enumerate}
        \item 长度 $\mathrm{L}$，单位：米 $\mathrm{m}$
        \item 时间 $\mathrm{T}$，单位：秒 $\mathrm{s}$
        \item 质量 $\mathrm{M}$，单位：千克 $\mathrm{kg}$
        \item 电流 $\mathrm{I}$，单位：安培 $\mathrm{A}$
        \item 热力学温度 $\varTheta$，单位：开尔文 $\mathrm{K}$
        \item 物质的量 $\mathrm{N}$，单位：摩尔 $\mathrm{mol}$
        \item 发光强度 $\mathrm{J}$，单位：坎德拉 $\mathrm{cd}$
    \end{enumerate}
    由基本量导出的物理量称为\textbf{导出量}，对应单位为导出单位。

    \subsubsection{量纲}
    物理量用基本量乘方乘积表达的形式；\textbf{量纲一致性原理}：物理方程等式两侧量纲必须完全相同，可用于校验公式正误。
    例：力的量纲 $\mathrm{MLT^{-2}}$；所有量纲指数为 $0$ 的量为无量纲量。

    \subsection{时间与空间}
    \begin{enumerate}
        \item 时空本质：客观存在，是运动物质的存在形式，二者通过物质运动相互关联，无法脱离物质独立存在。
        \item 时间定义：描述物质运动持续性与顺序性，1967年定义：铯-133原子基态两个超精细能级跃迁辐射9192631770个周期的时长为1秒。
        \item 长度定义：描述物质广延性，1983年定义：光在真空中 $1/299792458\ \mathrm{s}$ 内传播的路径长度为1米。
    \end{enumerate}

    \subsection{参考系和坐标系}
        \subsubsection{参考系}
        运动具有绝对性，观测描述具有相对性；被选定作为观测基准的物体称为\textbf{参考系}，所有物理量必须统一至同一参考系计算。

        \subsubsection{坐标系}
        在选定参考系中搭建的定量描述系统，常用：直角坐标系、自然坐标系、极坐标系、球坐标系、柱坐标系；坐标系选定等价于参考系选定，选取原则简化运算。

\section{质点运动的描述}
    \subsection{质点}
    理想化物理模型：当物体自身尺寸、形状对研究问题无影响时，抽象为集中全部质量的几何点。
    相对性示例：地球公转可视为质点，地球自转不可视为质点。

    \subsection{位置矢量与运动方程}
    \subsubsection{位置矢量（位矢）}
    参考系原点指向质点空间位置的矢量，记作 $\boldsymbol{r}$。
    直角坐标系分解：
    \[
    \boldsymbol{r}(t) = x(t)\boldsymbol{i} + y(t)\boldsymbol{j} + z(t)\boldsymbol{k}
    \]
    分量形式：
    \[
    \begin{cases}
    x = x(t)\\
    y = y(t)\\
    z = z(t)
    \end{cases}
    \]
    位矢大小：
    \[
    r = |\boldsymbol{r}| = \sqrt{x^2+y^2+z^2}
    \]
    方向由方向余弦确定：
    \[
    \cos\alpha = \frac{x}{r},\quad \cos\beta = \frac{y}{r},\quad \cos\gamma = \frac{z}{r}
    \]

    \subsubsection{运动方程与轨道方程}
    \begin{enumerate}
        \item 运动方程：$\boldsymbol{r}=\boldsymbol{r}(t)$，描述位矢随时间变化规律；
        \item 轨道方程：消去运动方程参数 $t$ 得到空间轨迹方程，分为直线运动、曲线运动。
    \end{enumerate}

    \subsection{位移与路程}
    \begin{enumerate}
        \item 位移矢量 $\Delta\boldsymbol{r}$：$\Delta t$ 内位矢的增量，矢量，方向由初位置指向末位置
        \[
        \Delta\boldsymbol{r} = \boldsymbol{r}_B - \boldsymbol{r}_A = \Delta x \boldsymbol{i} + \Delta y \boldsymbol{j} + \Delta z \boldsymbol{k}
        \]
        位移大小：
        \[
        |\Delta\boldsymbol{r}| = \sqrt{\Delta x^2+\Delta y^2+\Delta z^2}
        \]
        \item 路程 $\Delta s$：质点实际走过的轨迹长度，标量；
        \item 核心关系：一般 $|\Delta\boldsymbol{r}| \neq \Delta s$，仅单向直线运动或无穷小时间间隔内 $\mathrm{d}s = |\mathrm{d}\boldsymbol{r}|$。
    \end{enumerate}

    \subsection{速度矢量}
    \subsubsection{平均速度}
    $\Delta t$ 内位矢平均变化率，矢量：
    \[
    \bar{\boldsymbol{v}} = \frac{\Delta\boldsymbol{r}}{\Delta t}
    \]

    \subsubsection{瞬时速度（速度）}
    $\Delta t \to 0$ 时平均速度极限，沿轨迹切线方向：
    \[
    \boldsymbol{v} = \lim_{\Delta t \to 0}\frac{\Delta\boldsymbol{r}}{\Delta t} = \frac{\mathrm{d}\boldsymbol{r}}{\mathrm{d}t}
    \]
    直角坐标系分解：
    \[
    \boldsymbol{v} = \frac{\mathrm{d}x}{\mathrm{d}t}\boldsymbol{i} + \frac{\mathrm{d}y}{\mathrm{d}t}\boldsymbol{j} + \frac{\mathrm{d}z}{\mathrm{d}t}\boldsymbol{k} = v_x\boldsymbol{i}+v_y\boldsymbol{j}+v_z\boldsymbol{k}
    \]

    \subsubsection{速率}
    速度矢量的模，标量：
    \[
    v = |\boldsymbol{v}| = \sqrt{v_x^2+v_y^2+v_z^2}
    \]
    注意：平均速度大小 $\neq$ 平均速率。

\section{质点运动变化的描述}
    \subsection{加速度矢量}
    描述速度矢量（大小、方向）随时间变化快慢的物理量。
    \begin{enumerate}
        \item 平均加速度：
        \[
        \bar{\boldsymbol{a}} = \frac{\Delta\boldsymbol{v}}{\Delta t}
        \]
        \item 瞬时加速度：
        \[
        \boldsymbol{a} = \lim_{\Delta t \to 0}\frac{\Delta\boldsymbol{v}}{\Delta t} = \frac{\mathrm{d}\boldsymbol{v}}{\mathrm{d}t} = \frac{\mathrm{d}^2\boldsymbol{r}}{\mathrm{d}t^2}
        \]
    \end{enumerate}
    直角坐标系分解：
    \[
    \boldsymbol{a} = \frac{\mathrm{d}v_x}{\mathrm{d}t}\boldsymbol{i} + \frac{\mathrm{d}v_y}{\mathrm{d}t}\boldsymbol{j} + \frac{\mathrm{d}v_z}{\mathrm{d}t}\boldsymbol{k}
    \]
    加速度大小：
    \[
    a = \sqrt{a_x^2+a_y^2+a_z^2}
    \]

    \subsubsection{两类质点运动学核心问题}
    \begin{enumerate}
        \item 微分问题：已知运动方程 $\boldsymbol{r}(t)$，一阶导求速度，二阶导求加速度；
        \item 积分问题：已知加速度/速度与初始条件，积分逆运算求解速度、运动方程。
    \end{enumerate}
    \subsubsection{运动叠加原理}
    \textbf{运动叠加原理}：任意曲线运动可分解为三个正交方向独立直线运动的叠加。

    \subsection{自然坐标系中的速度和加速度}
    适用于已知轨迹的平面曲线运动，以切向单位矢量 $\boldsymbol{e}_\mathrm{t}$、法向单位矢量 $\boldsymbol{e}_\mathrm{n}$ 搭建坐标系。

    \subsubsection{自然坐标与速度}
    弧长坐标 $s=s(t)$，速率 $v=\dfrac{\mathrm{d}s}{\mathrm{d}t}$，速度矢量：
    \[
    \boldsymbol{v} = \frac{\mathrm{d}s}{\mathrm{d}t}\boldsymbol{e}_\mathrm{t} = v\boldsymbol{e}_\mathrm{t}
    \]

    \subsubsection{切向加速度与法向加速度}
    对速度求导分解：
    \[
    \boldsymbol{a} = \frac{\mathrm{d}v}{\mathrm{d}t}\boldsymbol{e}_\mathrm{t} + \frac{v^2}{\rho}\boldsymbol{e}_\mathrm{n}
    \]
    \begin{enumerate}
        \item 切向加速度：
        \[
        a_\mathrm{t} = \frac{\mathrm{d}v}{\mathrm{d}t}
        \]
        物理意义：改变速度的大小；
        \item 法向加速度（向心加速度）：
        \[
        a_\mathrm{n} = \frac{v^2}{\rho}
        \]
        物理意义：改变速度的方向，$\rho$ 为轨迹曲率半径；
    \end{enumerate}
    合加速度大小：
    \[
    a = \sqrt{a_\mathrm{t}^2 + a_\mathrm{n}^2}
    \]

    \subsubsection{由切向、法向加速度判定运动类型}
    \begin{enumerate}
        \item 若 $a_\mathrm{n} \equiv 0$：质点做直线运动；若 $a_\mathrm{n} \not\equiv 0$：质点做曲线运动；
        \item 若 $a_\mathrm{t} \equiv 0$：质点做匀速率运动；若 $a_\mathrm{t} \not\equiv 0$：质点做变速率运动。
    \end{enumerate}
    直线运动、匀速率曲线运动均为平面一般曲线运动的特例。

    \subsection{圆周运动的角量描述}
    \subsubsection{线量与角量定义}
    \begin{enumerate}
        \item \textbf{线量}：自然坐标系描述物理量，包括弧长 $s$、路程 $\Delta s$、速率 $v$、切向加速度 $a_\mathrm{t}$、法向加速度 $a_\mathrm{n}$；
        \item \textbf{角量}：极坐标系描述圆周运动物理量，包括角位置 $\theta$、角位移 $\Delta\theta$、角速度 $\omega$、角加速度 $\alpha$。
    \end{enumerate}
    规定：质点绕原点逆时针转动角量为正值，顺时针转动角量为负值。

    \subsubsection{角速度与角加速度定义式}
    角速度：
    \[
    \omega = \lim_{\Delta t \to 0}\frac{\Delta\theta}{\Delta t} = \frac{\mathrm{d}\theta}{\mathrm{d}t}
    \]
    角加速度：
    \[
    \alpha = \lim_{\Delta t \to 0}\frac{\Delta\omega}{\Delta t} = \frac{\mathrm{d}\omega}{\mathrm{d}t} = \frac{\mathrm{d}^2\theta}{\mathrm{d}t^2}
    \]
    SI单位：角速度 $\mathrm{rad\cdot s^{-1}}$，角加速度 $\mathrm{rad\cdot s^{-2}}$。

    \subsubsection{圆周运动角量与线量关系式}
    \[
    \begin{cases}
    v = R\omega \\
    a_\mathrm{t} = R\alpha \\
    a_\mathrm{n} = R\omega^2 = \dfrac{v^2}{R}
    \end{cases}
    \]

    \subsubsection{角速度与角加速度的矢量表达}
    无限小角位移满足矢量运算法则，瞬时角速度 $\boldsymbol{\omega}$、瞬时角加速度 $\boldsymbol{\alpha}$ 为矢量，方向由右手螺旋定则判定。
    线速度矢量叉乘公式：
    \[
    \boldsymbol{v} = \boldsymbol{\omega} \times \boldsymbol{r}
    \]
    加速度矢量叉乘公式：
    \[
    \boldsymbol{a} = \boldsymbol{\alpha} \times \boldsymbol{r} + \boldsymbol{\omega} \times \boldsymbol{v}
    \]
    式中第一项为切向加速度，第二项为法向加速度。

\section{相对运动}
    \subsection{参考系平动变换基本关系}
    设参考系 $S'$ 相对惯性系 $S$ 做平动，$S'$ 原点位矢为 $\boldsymbol{r}_0$，质点在 $S$ 系位矢 $\boldsymbol{r}$，在 $S'$ 系位矢 $\boldsymbol{r}'$，满足位矢变换：
    \[
    \boldsymbol{r} = \boldsymbol{r}' + \boldsymbol{r}_0
    \]
    位移变换关系：
    \[
    \Delta\boldsymbol{r} = \Delta\boldsymbol{r}' + \Delta\boldsymbol{r}_0
    \]

    \subsection{伽利略速度变换定理}
    对位矢公式对时间求导，令质点相对 $S$ 速度 $\boldsymbol{v}$，质点相对 $S'$ 速度 $\boldsymbol{v}'$，$S'$ 相对 $S$ 牵连速度 $\boldsymbol{u}$：
    \[
    \boldsymbol{v} = \boldsymbol{v}' + \boldsymbol{u}
    \]
    该式为\textbf{伽利略速度变换公式}，适用于物体运动速度远小于光速的宏观低速场景。

    \subsection{伽利略加速度变换定理}
    对速度变换式再次求导，$S'$ 相对 $S$ 牵连加速度 $\boldsymbol{a}_0$，质点在两参考系加速度 $\boldsymbol{a}$、$\boldsymbol{a}'$：
    \[
    \boldsymbol{a} = \boldsymbol{a}' + \boldsymbol{a}_0
    \]
    若 $S'$ 相对 $S$ 做匀速直线平动，牵连加速度 $\boldsymbol{a}_0 = \boldsymbol{0}$，则：
    \[
    \boldsymbol{a} = \boldsymbol{a}'
    \]
    即惯性系中所有匀速平动参考系内观测到的质点加速度完全相等。
% 4.质点动力学.tex 第二章 质点动力学
\chapter{第二章 质点动力学}
动力学研究物体间相互作用，以及相互作用引发的运动状态变化规律，是研究复杂物体运动的基础。本章以低速宏观运动为前提，围绕动量、角动量两大守恒量展开，并补充宇宙速度、黑洞物理模型与传送带动力学实例。

\section{动量 动量守恒定律}
    \subsection{动量与冲量}
    \subsubsection{质点动量定理}
    冲量定义：力对时间的累积效应 $\displaystyle \boldsymbol{I} = \int_{t_0}^{t}\boldsymbol{F}\mathrm{d}t$。
    质点动量定理：质点末态相对初态的动量增量等于所受合外力的冲量。
    \keyformula{\boldsymbol{I}=m\boldsymbol{v}-m\boldsymbol{v}_0}
    冲量 $\boldsymbol{I}$ 为矢量，方向与动量增量同向；单位 $\mathrm{N\cdot s}$，量纲 $\mathrm{MLT^{-1}}$。
    直角坐标系分量形式：
    \keyformula{
    \begin{cases}
    I_x = \displaystyle \int_{t_0}^{t}F_x\mathrm{d}t = mv_x - mv_{x0}\\
    I_y = \displaystyle \int_{t_0}^{t}F_y\mathrm{d}t = mv_y - mv_{y0}\\
    I_z = \displaystyle \int_{t_0}^{t}F_z\mathrm{d}t = mv_z - mv_{z0}
    \end{cases}
    }
    爆炸、碰撞瞬时冲力可用平均冲力简化：$\bar{\boldsymbol{F}}\mathrm{d}t = m\boldsymbol{v}-m\boldsymbol{v}_0$。

    \subsubsection{质点系动量定理}
    系统合外力等于总动量对时间的变化率：
    \[
    \boldsymbol{F}=\sum_i \boldsymbol{F}_i = \frac{\mathrm{d}\boldsymbol{p}}{\mathrm{d}t}
    \]
    积分形式质点系动量定理：
    \keyformula{\boldsymbol{I}=\int_{t_0}^{t}\boldsymbol{F}\mathrm{d}t = \boldsymbol{p}-\boldsymbol{p}_0=\Delta\boldsymbol{p}}
    系统内力矢量和为零，内力冲量无法改变系统总动量，仅能分配质点间动量。

    \subsubsection{动量守恒定律}
    系统合外力 $\boldsymbol{F}=\sum\boldsymbol{F}_i=\boldsymbol{0}$ 时，总动量为常矢量：
    \keyformula{\boldsymbol{p}=\boldsymbol{p}_0=\text{常矢量}}
    分量守恒条件：某一坐标轴外力分量为零，则该方向动量分量守恒
    \keyformula{
    \begin{cases}
    F_x=0,\quad p_x=\sum m_i v_{ix}=\text{常量}\\
    F_y=0,\quad p_y=\sum m_i v_{iy}=\text{常量}\\
    F_z=0,\quad p_z=\sum m_i v_{iz}=\text{常量}
    \end{cases}
    }
    适用范围：宏观物体、微观粒子（原子、分子、光子）普遍成立；碰撞爆炸问题可忽略重力等常规外力。

    \subsection{质心 质心运动定理}
    \subsubsection{质心位置矢量}
    离散质点系质心位矢（质量加权平均）：
    \keyformula{\boldsymbol{r}_C = \frac{\sum m_i \boldsymbol{r}_i}{m},\quad m=\sum m_i}
    直角坐标分量：
    \keyformula{
    \begin{cases}
    x_C = \dfrac{\sum m_i x_i}{m}\\[4pt]
    y_C = \dfrac{\sum m_i y_i}{m}\\[4pt]
    z_C = \dfrac{\sum m_i z_i}{m}
    \end{cases}
    }
    连续质量分布物体质心：
    \keyformula{\boldsymbol{r}_C = \frac{\int \boldsymbol{r}\mathrm{d}m}{m}}
    分量形式：
    \keyformula{
    \begin{cases}
    x_C = \dfrac{\int x\mathrm{d}m}{m}\\[4pt]
    y_C = \dfrac{\int y\mathrm{d}m}{m}\\[4pt]
    z_C = \dfrac{\int z\mathrm{d}m}{m}
    \end{cases}
    }

    \subsubsection{质心速度与加速度}
    \keyformula{\boldsymbol{v}_C = \frac{\mathrm{d}\boldsymbol{r}_C}{\mathrm{d}t} = \frac{\sum m_i \boldsymbol{v}_i}{m}}
    \keyformula{\boldsymbol{a}_C = \frac{\mathrm{d}\boldsymbol{v}_C}{\mathrm{d}t} = \frac{\sum m_i \boldsymbol{a}_i}{m}}

    \subsubsection{质心运动定理}
    系统总动量 $\boldsymbol{p}=m\boldsymbol{v}_C$，求导得：
    \keyformula{\boldsymbol{F}_{\text{外}} = \frac{\mathrm{d}\boldsymbol{p}}{\mathrm{d}t} = m\boldsymbol{a}_C}
    物理意义：系统合外力等于总质量乘以质心加速度；内力不改变质心运动状态，合外力为零时质心静止或匀速直线运动。

\section{角动量 角动量守恒定律}
    \subsection{质点角动量}
    质点对定点 $O$ 的角动量（动量矩）定义为位矢与动量叉乘：
    \keyformula{\boldsymbol{L} = \boldsymbol{r} \times \boldsymbol{p} = \boldsymbol{r} \times m\boldsymbol{v}}
    大小 $L=rmv\sin\varphi$，方向垂直 $\boldsymbol{r}$ 与 $\boldsymbol{v}$ 构成平面，右手螺旋定则判定；单位 $\mathrm{kg\cdot m^2\cdot s^{-1}}$，量纲 $\mathrm{ML^2T^{-1}}$。
    质点系总角动量：
    \keyformula{\boldsymbol{L}=\sum \boldsymbol{L}_i = \sum \boldsymbol{r}_i \times m_i \boldsymbol{v}_i}

    \subsection{力矩 角动量定理}
    \subsubsection{力矩定义}
    外力对定点力矩：
    \keyformula{\boldsymbol{M} = \boldsymbol{r} \times \boldsymbol{F}}
    大小 $M=rF\sin\varphi=Fd$，$d$ 为力臂；一对内力对同一参考点力矩矢量和恒为零 $\boldsymbol{M}_{\text{内}}\equiv \boldsymbol{0}$。
    质点系合外力矩：
    \keyformula{\boldsymbol{M}=\sum \boldsymbol{r}_i \times \boldsymbol{F}_i}

    \subsubsection{角动量定理}
    微分形式（质点/质点系通用）：
    \keyformula{\boldsymbol{M} = \frac{\mathrm{d}\boldsymbol{L}}{\mathrm{d}t}}
    文字表述：质点（质点系）所受对某一参考点的合外力矩，等于角动量对同一参考点随时间的变化率。

    积分形式（冲量矩定理）：
    \keyformula{\int_{t_1}^{t_2}\boldsymbol{M}\mathrm{d}t = \boldsymbol{L}_2 - \boldsymbol{L}_1}
    $\displaystyle \int_{t_1}^{t_2}\boldsymbol{M}\mathrm{d}t$ 称为**冲量矩**，是外力矩对时间的累积量；冲量矩等于系统角动量的增量。
    角动量单位：$\mathrm{m\cdot N\cdot s}$，量纲 $\mathrm{ML^2T^{-1}}$。

    \subsubsection{角动量守恒定律}
    若对某定点合外力矩 $\boldsymbol{M}=\boldsymbol{0}$，则系统对该点角动量 $\boldsymbol{L}=$ 常矢量，角动量守恒。
    \corepoint{有心力对力心的力矩恒为零，仅受有心力作用的系统角动量守恒；典型应用：开普勒行星运动第二定律（面积定律）}

    \subsubsubsection{例：开普勒第二定律推导}
    行星仅受太阳有心力，对太阳角动量守恒。行星角动量大小：
    \[
    L = rmv\sin\varphi = rm \left|\frac{\mathrm{d}\boldsymbol{r}}{\mathrm{d}t}\right|\sin\varphi
    \]
    极短时间 $\Delta t$ 内矢径扫过面积 $\displaystyle \Delta S = \frac12 r |\Delta \boldsymbol{r}|\sin\varphi$，取极限：
    \[
    L = 2m \lim_{\Delta t \to 0}\frac{\Delta S}{\Delta t} = 2m \frac{\mathrm{d}S}{\mathrm{d}t}
    \]
    \keyformula{\frac{\mathrm{d}S}{\mathrm{d}t} = \frac{L}{2m}}
    $L$ 为恒量，故面积速率 $\dfrac{\mathrm{d}S}{\mathrm{d}t}$ 不变，行星矢径在相等时间内扫过相等面积。

    \subsubsubsection{角动量守恒经典实例}
    1. 细绳拉小球绕细管转动：拉力为有心力，角动量守恒，半径缩小转速增大 $v_1 = \dfrac{r_2}{r_1}v_2$；
    2. 星系坍缩：万有引力为有心力，星云向内收缩时旋转角速度增大，形成盘状星系。

\section{变质量物体运动（火箭方程）}
    火箭动力学微分方程：
    \[
    m\frac{\mathrm{d}\boldsymbol{v}}{\mathrm{d}t} = \boldsymbol{F} + \boldsymbol{u}\frac{\mathrm{d}m}{\mathrm{d}t}
    \]
    $\boldsymbol{u}\dfrac{\mathrm{d}m}{\mathrm{d}t}$ 为火箭发动机推力，$\boldsymbol{u}$ 为喷出气体相对火箭速度。
    太空无外力时火箭末速度公式（齐奥尔科夫斯基公式）：
    \keyformula{v_m = v_0 + u\ln\frac{m_0}{m_{\text{火箭}}}}

\section{能量 能量守恒定律}
    \subsection{动能 功 动能定理}
    \subsubsection{动能}
    质点动能定义：
    \keyformula{E_\mathrm{k}=\frac12 mv^2}
    单位：焦耳（$\mathrm{J}$），量纲 $\mathrm{ML^2T^{-2}}$；
    质点系总动能为所有质点动能代数和：
    \keyformula{E_\mathrm{k}=\sum_i \frac12 m_i v_i^2}

    \subsubsection{功与功率}
    1. 元功定义：
    \keyformula{\mathrm{d}W = \boldsymbol{F}\cdot \mathrm{d}\boldsymbol{r}}
    2. 有限路径总功（线积分）：
    \keyformula{W=\int_{A}^{B}\boldsymbol{F}\cdot \mathrm{d}\boldsymbol{r}}
    直角坐标系分解：
    \[
    W=\int_L F_x\mathrm{d}x+F_y\mathrm{d}y+F_z\mathrm{d}z = W_x+W_y+W_z
    \tag{2-35}
    \]

    功四大核心性质：
    \begin{enumerate}[label=(\alph*)]
        \item 标量，正负由力与位移夹角决定：$\alpha<\pi/2$ 正功，$\alpha>\pi/2$ 负功，$\alpha=\pi/2$ 不做功；合外力做功等于各分力做功代数和；
        \item 过程量：功依赖完整运动路径，某一时刻无“功”的概念；
        \item 相对量：位移与参考系绑定，不同参考系下同一力做功不同；
        \item 一对相互作用力总功由相对位移决定，$\mathrm{d}W=\boldsymbol{F}_1\cdot \mathrm{d}\boldsymbol{r}_{12}$，无相对位移时内力总功为0。
    \end{enumerate}

    瞬时功率定义：
    \keyformula{P=\lim_{\Delta t\to0}\frac{\Delta W}{\Delta t}=\frac{\mathrm{d}W}{\mathrm{d}t}=\boldsymbol{F}\cdot \boldsymbol{v}}
    单位：瓦特（$\mathrm{W}$）。

    \subsubsection{质点动能定理}
    对元功公式沿路径积分，得动能定理：
    \keyformula{W=\int_{A}^{B}\boldsymbol{F}\cdot \mathrm{d}\boldsymbol{r} = \Delta E_\mathrm{k} = \frac12 mv_B^2 - \frac12 mv_A^2}
    \corepoint{动能定理核心：合外力对质点做的总功，等于质点动能的增量；功是过程量，动能是状态量，只与初末速率有关}
    适用范围：惯性参考系，处理速率–位置关系问题比牛顿第二定律更简便。

    \subsubsection{传送带滑动摩擦动力学（动能定理应用）}
    质量$m$物体轻放在速率$v$的水平传送带上，动摩擦因数$\mu$，滑动摩擦力$F_f=\mu mg$。
    1. 加速滑动时间（动量定理）
    \[
    -F_f \cdot \Delta t = 0-mv \implies \Delta t = \frac{v}{\mu g}
    \]
    2. 物体对地滑动位移（动能定理）
    \[
    -F_f \cdot s = 0-\frac12 mv^2 \implies s = \frac{v^2}{2\mu g}
    \]
    3. 摩擦耗散机械能
    传送带对地位移$s'=v\Delta t=\dfrac{v^2}{\mu g}$，相对位移$s'-s$，摩擦生热：
    \[
    W' = \mu mg\left(\frac{v^2}{\mu g}-\frac{v^2}{2\mu g}\right) = \frac12 mv^2
    \]

    \subsection{保守力 势能 功能原理}
    \subsubsection{保守力判定}
    保守力：做功仅由初末位置决定、与运动路径无关的力，满足环路积分 $\displaystyle \oint_L \boldsymbol{F}_\text{保}\cdot \mathrm{d}\boldsymbol{r}=0$；
    典型保守力：重力、弹性力、万有引力、静电力；
    非保守力（耗散力）：摩擦力、粘滞阻力、爆炸推力，做功与路径相关。

    \subsubsection{三类保守力的功与势能}
    1. 重力的功与重力势能
    \[
    W_G = -mg(y_b-y_a)
    \tag{2-41}
    \]
    取地面为势能零点，高度 $h$ 处重力势能：
    \keyformula{E_{\mathrm{p}\text{重}}=mgh}

    2. 弹性力的功与弹性势能
    \[
    W_\text{弹}= -\left(\frac12 kx_b^2 - \frac12 kx_a^2\right)
    \tag{2-42}
    \]
    弹簧原长为势能零点，形变量 $x$ 弹性势能：
    \keyformula{E_{\mathrm{p}\text{弹}}=\frac12 kx^2}

    3. 万有引力的功与引力势能
    \[
    W_\text{引} = -Gm_1m_2\left(\frac{1}{r_b}-\frac{1}{r_a}\right)
    \tag{2-43}
    \]
    取无穷远为势能零点，间距 $r$ 引力势能：
    \keyformula{E_{\mathrm{p}\text{引}}=-G\frac{m_1m_2}{r}}

    \subsubsection{保守力与势能梯度关系}
    保守力做功等于势能增量的负值：
    \[
    W_\text{保} = -\Delta E_\mathrm{p},\quad \mathrm{d}W_\text{保}=-\mathrm{d}E_\mathrm{p}
    \]
    \keyformula{\boldsymbol{F}_\text{保}=-\nabla E_\mathrm{p} = -\left(\frac{\partial E_\mathrm{p}}{\partial x}\boldsymbol{i}+\frac{\partial E_\mathrm{p}}{\partial y}\boldsymbol{j}+\frac{\partial E_\mathrm{p}}{\partial z}\boldsymbol{k}\right)}
    物理意义：保守力等于对应势能函数的负梯度；势能属于相互作用系统共有。

    \subsubsection{质点系功能原理}
    将内力拆分为保守内力、非保守内力，推导得功能原理：
    \keyformula{W_\text{外} + W_\text{非保内} = E - E_0}
    其中机械能 $E=E_\mathrm{k}+E_\mathrm{p}$；文字表述：系统外力做功与非保守内力做功之和，等于系统机械能的增量。

    \subsection{机械能守恒定律与能量守恒定律}
    \subsubsection{机械能守恒定律}
    守恒条件：$W_\text{外}=0$ 且 $W_\text{非保内}=0$，只有保守内力做功：
    \keyformula{E=E_\mathrm{k}+E_\mathrm{p}=\text{常量}}
    \corepoint{机械能守恒仅适用于惯性系；碰撞分类：弹性碰撞（机械能守恒）、完全非弹性碰撞、非弹性碰撞（后两者机械能损耗）}

    \subsubsection{普遍能量守恒定律}
    孤立系统（与外界无能量交换）内，一切形式能量仅相互转化，总能量保持恒定；能量既不能凭空产生，也不能凭空消失，是自然界普适基本定律。

    \subsubsection{经典应用：宇宙速度与黑洞模型}
    1. 第二宇宙速度（脱离地球逃逸速度）
    以无穷远为引力势能零点，脱离地球引力末速度 $v_\infty=0$，机械能守恒：
    \[
    \frac12 mv^2 - G\frac{mM}{R}=0
    \]
    代入地面重力关系 $G\dfrac{M}{R^2}=g$，得第二宇宙速度：
    \keyformula{v=\sqrt{2gR}}
    推广到任意天体，逃逸速度通式：
    \[
    v=\sqrt{\frac{2Gm_{\text{天体}}}{R}}
    \]
    2. 黑洞物理简述
    若天体密度极大、半径极小，逃逸速度可达到光速；此时天体发射的光子无法挣脱引力束缚，外部无法观测该天体，此类天体定义为黑洞（black hole），是天体物理前沿研究课题。

\endinput
\chapter{第三章 刚体力学基础}

\section{刚体运动的基本形式}
    \subsection{平动和转动}
    刚体可以看成是由无数个连续分布的质点所组成的质点系，每个质点称为刚体的一个质量元。刚体在运动中各部分相对位置均保持不变，因此只要知道刚体上某一点的运动情况，则刚体上任一点乃至刚体运动的全貌便可确定。

    刚体的运动可以分为平动和转动，刚体任何复杂的运动均可看成是二者的叠加。刚体上任一给定直线在运动中空间方向始终不变而保持平行，这样的运动称为平动。刚体平动时，刚体上各点的运动轨迹均相同，同一时刻刚体上各点均具有相同的速度和加速度，刚体平动时完全可以简化为质点处理，通常使用刚体质心的运动代表整个刚体的平动。

    如果刚体上所有质点都绕同一直线作圆周运动，这种运动称为刚体的转动，这条直线称为转轴。若整个转轴相对参考系静止，这种转动称为定轴转动；假如转轴上只有一点相对参考系静止，而转轴方向却不断变动，我们称这种转动为定点转动。刚体的任何复杂运动均可看成是平动和转动的叠加，刚体的定轴转动及其规律是本章的重点内容。

    \subsection{描述刚体转动的物理量}
    刚体定轴转动时，轴外刚体上的各个质点，尽管到轴的距离不同，相同的时间内转过的线位移也不同，但转过的角位移则相同。因此可以在刚体上任意选定一点，通过研究该点绕定轴的转动，以此来描述刚体的定轴转动。

    刚体上每个质点的转动平面都是与转轴相垂直的平面。通常选取与轴相固联的惯性系为参考系，相对参考系的任一固定方向 $Ox$ 均可作为参考方向，质点对转轴交点 $O$ 的位矢 $\boldsymbol{r}$ 与该 $Ox$ 方向的夹角 $\theta$ 即为质点的角坐标，刚体转动时角速度为
    \[
    \omega=\frac{\mathrm{d}\theta}{\mathrm{d}t}
    \]
    角速度方向沿转轴，遵从右手螺旋定则。当角速度变化时，角加速度大小为
    \[
    \alpha=\frac{\mathrm{d}\omega}{\mathrm{d}t}=\frac{\mathrm{d}^2\theta}{\mathrm{d}t^2}
    \]
    $\omega$ 变大时 $\alpha$ 与 $\omega$ 同向，$\omega$ 变小时 $\alpha$ 与 $\omega$ 反方向。

    在刚体定轴转动中，角速度、角加速度方向只有沿转轴的两个方向，所以计算中常作标量处理。相应质点的线速度 $\boldsymbol{v}$ 和加速度 $\boldsymbol{a}$ 与角量的关系：
    \[
    \begin{cases}
    \boldsymbol{v}=\boldsymbol{\omega}\times\boldsymbol{r},\\
    \boldsymbol{a}=\boldsymbol{\alpha}\times\boldsymbol{r}+\boldsymbol{\omega}\times\boldsymbol{v},
    \end{cases}
    \]
    其中，$\boldsymbol{\alpha}\times\boldsymbol{r}$ 是质点的切向加速度，大小 $a_\mathrm{t}=r\alpha$；$\boldsymbol{\omega}\times\boldsymbol{v}$ 是质点的法向加速度，大小 $a_\mathrm{n}=r\omega^2$。

\section{定轴转动刚体的角动量 转动惯量}
    \subsection{刚体对定轴的角动量}
    当刚体绕定轴作定轴转动时，刚体上各质元某一瞬时均以相同的角速度 $\boldsymbol{\omega}$ 绕该定轴作平面圆周运动。设刚体上某一质元 $\Delta m_i$ 距轴的距离为 $r_i$，其对该轴的角动量为 $L_i=\Delta m_i r_i^2 \omega$，整个刚体对该轴的角动量 $L_z$ 应等于刚体上所有质元对轴角动量的总和，即
    \[
    L_z = \sum_{i} L_i = \left(\sum_{i} \Delta m_i r_i^2\right)\omega
    \]

    \subsection{转动惯量}
    \textbf{转动惯量定义}：质点系内每个质点的质量与该质点到转轴距离平方之积的总和，用 $J$ 表示，即
    \[
    J = \sum_{i} \Delta m_i r_i^2
    \]
    对质量连续分布的刚体，取微元形式积分：
    \[
    J = \int_{m} r^2 \mathrm{d}m
    \]
    在国际单位制中，转动惯量单位为 $\mathrm{kg\cdot m^2}$，量纲为 $\mathrm{ML^2T^0}$。

    转动惯量 $J$ 的大小与刚体的总质量及质量相对于给定轴的分布有关；在总质量一定的情况下，质量分布离轴愈远，转动惯量愈大；同时转动惯量还与转轴的位置有关，同一刚体，转轴位置不同，转动惯量也就不同。

    引入转动惯量概念后，定轴转动刚体的角动量可以表示为
    \[
    L_z = J \omega
    \]
    刚体对转轴的角动量 $L_z$ 的方向与角速度 $\boldsymbol{\omega}$ 的方向一致，计算中通常设定正方向后用正负来表示 $L_z$ 的方向。

    \subsubsection{回转半径}
    设物体的总质量 $m=\sum \Delta m_i$，刚体对给定轴的转动惯量为 $J$，定义物体对该转轴的回转半径 $r_G$ 为
    \[
    r_G=\sqrt{\frac{J}{m}}
    \]

    \subsubsection{平行轴定理}
    \textbf{平行轴定理}：若质量为$m$的刚体对过其质心$C$的转轴的转动惯量为$J_C$，则刚体对平行于该轴、和它相距为$d$的另一平行转轴的转动惯量：
    \[
    J = J_C + md^2
    \]

    \subsubsection{正交轴定理}
    薄板形刚体对于板面内的两条正交轴的转动惯量之和等于这个物体对过该两轴交点并垂直于板面的那条转轴的转动惯量，即
    \[
    J_z = J_y + J_x
    \]

\section{刚体定轴转动定律}
    \subsection{刚体角动量的时间变化率 对转轴的力矩}
    质点系角动量对时间的变化率等于系统所受的合外力矩，即 $\dfrac{\mathrm{d}\boldsymbol{L}}{\mathrm{d}t}=\boldsymbol{M}$。

    力作用于定轴转动的刚体上时，外力 $\boldsymbol{F}$ 可以分解为平行于转轴的力 $\boldsymbol{F}_\parallel$ 和垂直于转轴（在转动平面内）的力 $\boldsymbol{F}_\perp$。$\boldsymbol{F}_\parallel$ 的力矩对转动不起作用，只有在转动平面里的 $\boldsymbol{F}_\perp$ 的力矩才对刚体的转动运动有作用。定轴转动刚体角动量的时间变化率与外力对转轴的力矩的矢量关系：
    \[
    \frac{\mathrm{d}\boldsymbol{L}}{\mathrm{d}t}=\boldsymbol{M}_z=\boldsymbol{r}\times\boldsymbol{F}_\perp
    \]
    力矩表征着力对物体转动的作用，它是力的大小、方向及作用点诸因素组成的一个物理量。

    \subsection{刚体定轴转动定律}
    \textbf{刚体定轴转动定律}：刚体所受的对于某定轴的合外力矩等于刚体对此定轴的转动惯量与刚体在此合外力矩作用下所获得的角加速度的乘积。
    \[
    M_z = J \alpha
    \]
    该式是合外力矩对刚体的瞬时作用规律，式中各量均需是同一时刻对同一刚体、同一转轴而言。

\section{刚体的角动量定理和角动量守恒定律}
    \subsection{刚体的角动量定理}
    刚体可以看作由质元组成的质点系，角动量定理微分形式和积分形式：
    \[
    \boldsymbol{M}=\frac{\mathrm{d}\boldsymbol{L}}{\mathrm{d}t},\qquad \int_{t_1}^{t_2} \boldsymbol{M}\mathrm{d}t=\boldsymbol{L}_2-\boldsymbol{L}_1
    \]
    对于定轴转动刚体，各矢量均可沿轴用标量表示，定轴角动量定理：
    \[
    \int_{t_1}^{t_2} M_z \mathrm{d}t = L_{z2}-L_{z1}
    \]
    该式表明：在一段时间内作用在刚体上的外力的冲量矩等于刚体在该段时间内的角动量的增量。

    \subsection{刚体的角动量守恒定律}
    若刚体所受合外力矩 $\boldsymbol{M}=0$，则有
    \[
    J\boldsymbol{\omega}=\text{常矢量}
    \]
    说明：若刚体所受合外力矩为零，则刚体的角动量保持不变，称为刚体的角动量守恒定律。

    刚体角动量守恒常用情况：
    \begin{enumerate}
        \item 对定轴转动的刚体，转动过程中转动惯量 $J$ 保持不变，合外力矩为零时，刚体将以恒定的角速度 $\omega$ 绕定轴转动。轮船、飞机、火箭上用作导航定向的回转仪就是利用这一原理制成的。
        \item 对定轴转动的非刚体，物体上各质元相对转轴距离可变，转动惯量 $J$ 也是可变的。当转动系统所受合外力矩 $M=0$ 时，$J\boldsymbol{\omega}=$常矢量，$J$ 和 $\boldsymbol{\omega}$ 成反比。人站在可绕竖直转轴转动的转台上，两臂平伸，两手各握重物，转动起来后收拢双臂，人和转台转速加快；伸出双臂，转速又变慢。滑冰运动员、跳水运动员、芭蕾演员，经常运用角动量守恒的原理完成优美动作。
        \item 物体系统的角动量守恒，当几个相互关联的质点和刚体组成系统时，各物体对同一转轴的角动量分别是 $J_1\boldsymbol{\omega}_1$，$J_2\boldsymbol{\omega}_2$，$\dots$，$\boldsymbol{r}_j\times m_j \boldsymbol{v}_j$，若整个系统对同一转轴所受的合外力矩为零，则整个物体系统对该转轴的总角动量守恒，即
        \[
        \sum_i J_i \boldsymbol{\omega}_i + \sum_j (\boldsymbol{r}_j \times m_j \boldsymbol{v}_j) = \text{常矢量}.
        \]
    \end{enumerate}
    角动量守恒定律与动量守恒定律和能量守恒定律一样，均是自然界中的普遍规律。在微观领域内，常应用角动量守恒来推算粒子衰变产物的运动状态。

\section{进动}
    陀螺是常见的玩具，当它不旋转时，由于受到重力矩的作用，会立即倾倒在地。但当陀螺绕自身对称轴高速旋转（通常称自旋）时，尽管同样也受到重力矩的作用，却不会倒下来。这是刚体转动轴不固定时的常见实例。我们把陀螺高速自转的同时，其自身对称轴还将绕竖直轴作回旋运动的现象称为进动。

    用角动量定理可以解释陀螺的进动。设陀螺的质量为 $m$，对自身轴的转动惯量为 $J$。若陀螺自转角速度为 $\boldsymbol{\omega}$，则陀螺绕自身轴的自旋角动量为
    \[
    \boldsymbol{L}=J\boldsymbol{\omega},
    \]
    式中，$\boldsymbol{L}$ 的方向沿自身的对称轴。以陀螺与地面的接触点 $O$ 为参考点，陀螺所受重力矩为
    \[
    \boldsymbol{M}=\boldsymbol{r}_c \times m\boldsymbol{g},
    \]
    $\boldsymbol{M}$ 的方向垂直纸面向里，与陀螺的自旋角动量 $\boldsymbol{L}$ 垂直。所以，重力矩不能改变 $\boldsymbol{L}$ 的大小，而只能改变其方向。当重力矩作用于陀螺 $\mathrm{d}t$ 时间后，陀螺自旋角动量 $\boldsymbol{L}$ 的增量为
    \[
    \mathrm{d}\boldsymbol{L}=\boldsymbol{M}\mathrm{d}t,
    \]
    $\mathrm{d}\boldsymbol{L}$ 的方向与 $\boldsymbol{M}$ 方向平行，即自旋角动量 $\boldsymbol{L}$ 的方向将水平地转向 $\boldsymbol{L}+\mathrm{d}\boldsymbol{L}$ 方向，并不沿竖直方向向下倾斜。于是，自旋角动量 $\boldsymbol{L}$ 在水平面内持续偏转而形成绕竖直轴的进动，表现为沿一个圆锥面转动。这一圆锥面的轴线是竖直的，锥顶就在陀螺尖顶与地面的接触点 $O$。由矢量图可以看出，在 $\mathrm{d}t$ 时间内，陀螺自转轴转过的角度为
    \[
    \mathrm{d}\varphi=\frac{|\mathrm{d}\boldsymbol{L}|}{L\sin\theta},
    \]
    相应的角速度称为进动角速度 $\boldsymbol{\Omega}$，其大小为
    \[
    \Omega=\frac{\mathrm{d}\varphi}{\mathrm{d}t}=\frac{|\mathrm{d}\boldsymbol{L}|}{L\sin\theta \mathrm{d}t}=\frac{M}{L\sin\theta}.
    \]
    式中，$\theta$ 为陀螺的自身轴与圆锥的轴线之间的夹角。上式说明，进动的角速度与外力矩 $M$ 成正比，而与自旋角动量成反比，亦即与自旋转角速度 $\boldsymbol{\omega}$ 成反比。

    使车轮绕水平放置的自身轴旋转起来，松手后，车轮的轴将在水平面内以竖直轴顶端的凹槽为中心旋转起来。我们把表现这种现象的具有自身对称轴的物体系统叫作回转仪，因而进动运动又叫回转运动，可以证明，其绕竖直轴进动的角速度仍具有相同的结果。

    进动效应在工程实践中有着广泛的应用。例如，炮弹在飞行过程中，要受到空气阻力，阻力的方向总是与炮弹质心的速度方向相反，但其力线不一定通过质心。阻力对质心的力矩就会使炮弹在空中翻转。为避免这种事故，就在炮膛内壁刻出螺旋线（或称来复线），使炮弹出膛后，还能绕自身轴高速旋转，这样，飞行中的空气阻力矩将不再使它翻转，而只是使它绕着质心前进方向进动。它的轴线始终只与前进方向有不大的偏离，而弹头总是大致指向前方了。

    地球本身就是一个很大的回转仪，因为地球有自转，又受到太阳及其他星体的引力，因而地球在运动中要产生进动，在微观世界中，进动现象也十分普遍。例如，物质中绕核运动的电子处在外磁场中时，也会产生以外磁场方向为轴线的进动。

    应当指出，进动角速度公式只是一个近似关系，因为在推导中作了一个简化处理：即把陀螺绕自身轴的角动量作为陀螺的总角动量，而实际上陀螺的总角动量应该是它的自旋角动量和进动角动量的矢量和。只有在自旋角速度 $\omega$ 远大于进动角速度 $\Omega$ 时，这样的简化才会有效。即只有 $\omega\gg\Omega$ 时，公式才成立。在一般情况下，陀螺（或车轮）的自身轴还会周期性地摆动（称章动）。

\section{刚体定轴转动的动能定理和机械能守恒定律}
    \subsection{刚体转动动能和势能}
    设刚体作定轴转动，在某一时刻角速度的大小为 $\omega$，刚体上任一质元 $\Delta m_i$ 距转轴 $r_i$，则它的动能为
    \[
    \Delta E_{\mathrm{k}_i}=\frac12 \Delta m_i v_i^2=\frac12 \Delta m_i r_i^2 \omega^2,
    \]
    整个刚体的转动动能为刚体所有质元动能之和，即
    \[
    E_\mathrm{k}=\sum_i \Delta E_{\mathrm{k}_i}=\frac12 \left(\sum_i \Delta m_i r_i^2\right)\omega^2.
    \]
    由转动惯量定义，$\sum_i \Delta m_i r_i^2 = J$，所以
    \[
    E_\mathrm{k}=\dfrac{1}{2}J\omega^2
    \]

    同样的，刚体重力势能为刚体所有质元重力势能之和。若任一质元 $\Delta m_i$，其相对势能零点的高度为 $h_i$，则刚体重力势能 $E_\mathrm{p}$ 为
    \[
    E_\mathrm{p}=\sum_i \Delta m_i g h_i,
    \]
    又可知质心相对势能零点的高度为
    \[
    h_C=\frac{\sum_i \Delta m_i h_i}{\sum_i \Delta m_i},
    \]
    考虑到刚体总质量 $m=\sum \Delta m_i$，所以刚体的重力势能为
    \[
    E_\mathrm{p}=m g h_C,
    \]
    即刚体的重力势能相当于刚体的质量 $m$ 集中在质心 $C$ 处的质点的重力势能，$h_C$ 为刚体质心高度。

    \subsection{力矩的功}
    外力对刚体做功等于作用于刚体的各个外力对刚体做功的总和。设刚体在外力作用下有一个元角位移 $\mathrm{d}\theta=\omega \cdot \mathrm{d}t$，外力 $\boldsymbol{F}_i$ 所做的元功 $\mathrm{d}W_i=\boldsymbol{F}_i \cdot \mathrm{d}\boldsymbol{l}_i$，其中 $\mathrm{d}\boldsymbol{l}_i$ 是 $\boldsymbol{F}_i$ 作用点在该点转动平面上的元位移，$\mathrm{d}\boldsymbol{l}_i=\boldsymbol{v}_i \mathrm{d}t$，而速度 $\boldsymbol{v}_i=\boldsymbol{\omega}\times\boldsymbol{r}_i$，所以
    \[
    \mathrm{d}W_i=\boldsymbol{F}_i \cdot (\boldsymbol{\omega}\times\boldsymbol{r}_i)\mathrm{d}t,
    \]
    又由矢量混积的轮换性 $\boldsymbol{a}\cdot(\boldsymbol{b}\times\boldsymbol{c})=\boldsymbol{b}\cdot(\boldsymbol{c}\times\boldsymbol{a})=\boldsymbol{c}\cdot(\boldsymbol{a}\times\boldsymbol{b})$，上式可记为
    \[
    \mathrm{d}W_i=\boldsymbol{\omega}\cdot(\boldsymbol{r}_i\times\boldsymbol{F}_i)\mathrm{d}t,
    \]
    显然，所有外力元功的总和为
    \[
    \mathrm{d}W=\sum_i \boldsymbol{\omega}\cdot(\boldsymbol{r}_i\times\boldsymbol{F}_i)\mathrm{d}t=\sum_i \boldsymbol{M}_i \cdot \boldsymbol{\omega}\mathrm{d}t=\boldsymbol{M}_\text{外}\cdot \mathrm{d}\boldsymbol{\theta},
    \]
    上式表明，力对刚体做功可以用力矩大小与角位移的乘积来计算，所以也叫力矩做功。当角位置从 $\theta_1\rightarrow\theta_2$，则力矩做功为
    \[
    W=\int \mathrm{d}W=\int_{\theta_1}^{\theta_2} \boldsymbol{M}_\text{外} \cdot \mathrm{d}\boldsymbol{\theta},
    \]
    由功率的定义，力矩的功率可表示为
    \[
    P=\frac{\mathrm{d}W}{\mathrm{d}t}=\frac{\boldsymbol{M}_\text{外}\cdot \mathrm{d}\boldsymbol{\theta}}{\mathrm{d}t}=\boldsymbol{M}_\text{外}\cdot \boldsymbol{\omega}.
    \]

    \subsection{刚体定轴转动的动能定理}
    质点动能定理可由牛顿定律导出，同样刚体定轴转动的动能定理也可由转动定轴定律导出。由力矩元功表达式
    \[
    \mathrm{d}W=M\mathrm{d}\theta=J\alpha \mathrm{d}\theta=J\frac{\mathrm{d}\omega}{\mathrm{d}t}\mathrm{d}\theta=J\mathrm{d}\omega \frac{\mathrm{d}\theta}{\mathrm{d}t}=J\omega \mathrm{d}\omega.
    \]
    在外力矩作用下，当刚体从 $t_1$ 时刻的 $\theta_1$ 和 $\omega_1$ 变化到 $t_2$ 时刻的 $\theta_2$ 和 $\omega_2$ 时，对上式积分得
    \[
    W=\int_{\theta_1}^{\theta_2} M \mathrm{d}\theta=\int_{\omega_1}^{\omega_2} J \omega \mathrm{d}\omega=\frac12 J \omega_2^2-\frac12 J \omega_1^2
    \]
    上式称为刚体定轴转动的动能定理。它表明刚体转动动能的增量等于合外力矩对刚体所做的功。

    \subsection{刚体定轴转动的机械能守恒定律}
    如果刚体定轴转动中除受外力矩外，尚受摩擦力矩和重力矩作用，按动能定理则有
    \[
    W=\int_{\theta_1}^{\theta_2} \left(\boldsymbol{M}_\text{外}+\boldsymbol{M}_\text{非保内}+\boldsymbol{M}_\text{重}\right)\mathrm{d}\theta=\left(\frac12 J \omega_2^2+m g h_{C2}\right)-\left(\frac12 J \omega_1^2+m g h_{C1}\right).
    \]
    上式即为重力场中刚体定轴转动的功能原理。

    如果刚体只受保守力矩作用，则刚体机械能守恒。即
    \[
    \frac12 J\omega^2 + m g h_C = \text{常量}.
    \]
    对于包含刚体的力学系统，只有保守内力做功，外力矩与非保守内力矩做功均为零时，系统的机械能（转动动能、平动动能、重力势能之和）保持守恒。

\section{刚体的平面平行运动}
    \subsection{刚体平面平行运动的描述}
    如果刚体上任一质元的运动轨迹都平行于某一固定平面，这种运动称为刚体的平面平行运动。刚体的平面平行运动的特点是刚体上每一质元的运动轨迹都是平面曲线，且各平面互相平行；刚体在运动中转轴始终保持平行且垂直于某一固定平面。

    刚体作平面平行运动时，通常的研究方法是将其运动分解为整个刚体随质心的平动和绕过质心并垂直于运动平面的转轴的定轴的转动。

    刚体作平面平行运动时，平动部分需用其质心速度 $\boldsymbol{v}_C$ 和质心加速度 $\boldsymbol{a}_C$ 来描述。而刚体绕质心的转动可用角速度 $\boldsymbol{\omega}$ 和角加速度 $\boldsymbol{\alpha}$ 来描述。这样，刚体上任一点的速度和加速度均可按运动叠加原理求出。据此，刚体上任一点（距质心位置为 $\boldsymbol{r}$）的速度可表示为
    \[
    \boldsymbol{v}=\boldsymbol{v}_C+\boldsymbol{\omega}\times\boldsymbol{r},
    \]
    任一点的加速度可表示为
    \[
    \boldsymbol{a}=\boldsymbol{a}_C+\boldsymbol{a}_\text{绕C},
    \]
    式中，$\boldsymbol{\omega}$ 和 $\boldsymbol{a}_\text{绕C}$ 均为刚体上该点绕质心转动时，在该时刻所具有的角速度和线加速度。

    \subsection{刚体平面平行运动的动力学规律}
    对刚体平面平行运动的动力学分析，同样可把运动分解为刚体随质心的平动和绕质心的转动来分析。

    刚体随质心的平动规律遵从质心运动定理，即
    \[
    \sum_i \boldsymbol{F}_i = m \frac{\mathrm{d}^2 \boldsymbol{r}_C}{\mathrm{d}t^2}=m \frac{\mathrm{d}\boldsymbol{v}_C}{\mathrm{d}t}=m \boldsymbol{a}_C.
    \]
    式中，$m$ 为刚体总质量，$\boldsymbol{a}_C$ 为质心的加速度，$\sum_i \boldsymbol{F}_i$ 为整个刚体受到的外力的合矢量（各外力 $\boldsymbol{F}_i$ 不一定通过质心）。刚体作平面平行运动时，质心运动定理的分量式为
    \[
    \begin{cases}
    \sum_i F_{xi}=m a_{Cx},\\
    \sum_i F_{yi}=m a_{Cy}.
    \end{cases}
    \]

    理论和实验均证明，刚体绕通过质心的转轴的转动遵从定轴转动定律，于是有
    \[
    M_C = J_C \alpha
    \]
    式中，$M_C$ 是外力对通过质心的转轴的合外力矩；$J_C$、$\alpha$ 是刚体对通过质心的转轴的转动惯量和角加速度。因此，质心平动公式和绕质心转动公式是求解刚体平面平行运动的动力学基本方程组。在刚体既有随质心平动又有绕质心转动的情况下，它的动能也包括随质心的平动和绕质心的转动两部分，即
    \[
    E_\mathrm{k}=\frac12 m v_C^2+\frac12 J_C \omega^2.
    \]

    \subsection{纯滚动}
    若刚体在滚动过程中刚体与接触面之间无相对滑动，这种运动称为纯滚动。

    当圆柱体或圆球作纯滚动时，质心的平动与绕质心轴的转动这两种运动虽然是独立的，但描述质心平动与绕质心转动的物理量之间却有一定的对应关系。圆柱体（或圆球）每滚过一周，其中心轴前进的距离必等于圆柱体横截面的周长。于是得出中心轴前进的距离 $x_C$ 与圆柱体绕中心轴转动的角度 $\theta$ 的关系为
    \[
    x_C=R\theta,
    \]
    式中，$R$ 为圆柱体横截面的半径。对上式求导得质心的前进速度 $v_C$ 与绕轴转动的角速度 $\omega$ 之间关系为
    \[
    v_C=\frac{\mathrm{d}x_C}{\mathrm{d}t}=R\frac{\mathrm{d}\theta}{\mathrm{d}t}=R\omega,
    \]
    再对上式求导，可得质心的加速度 $a_C$ 与转动角加速度间的关系为
    \[
    a_C=\frac{\mathrm{d}v_C}{\mathrm{d}t}=R\frac{\mathrm{d}\omega}{\mathrm{d}t}=R\alpha.
    \]
    我们将质心平动与绕质心轴转动的上述三个关系式称为纯滚动条件（求解问题时只需其中一个），即
    \[
    \begin{cases}
    x_C=R\theta,\\
    v_C=R\omega,\\
    a_C=R\alpha.
    \end{cases}
    \]
    通常应用质心运动定理、绕质心转动定律和纯滚动条件，即可求解作纯滚动的刚体平面平行运动问题。

\chapter{第四章 流体力学简介}

\section{流体运动的描述}
    \subsection{速度场和定常流动}
    定常流动：流体空间各点流速不随时间发生变化。
    \subsection{流线和流管}
    流线为流体微团运动轨迹，流管由闭合流线围成，理想流体定常流动中流管内流体质量守恒。

\section{连续性方程和伯努利方程}
    \subsection{理想流体}
    不可压缩、无黏滞阻力的理想化流体模型。
    \subsection{理想流体的连续性方程}
    \textbf{连续性方程}（定常不可压缩流体）：
    \[
    S_1 v_1 = S_2 v_2
    \]
    物理意义：同一流管不同截面体积流量相等。
    \subsection{理想流体的伯努利方程}
    \textbf{伯努利方程}（理想流体定常流动）：
    \[
    p_1+\frac{1}{2}\rho v_1^2+\rho g h_1 = p_2+\frac{1}{2}\rho v_2^2+\rho g h_2
    \]
    本质为流体单位体积机械能守恒。
\chapter{第五章 狭义相对论}

\section{狭义相对论基本原理}
    \subsection{伽利略相对性原理}
    \textbf{伽利略相对性原理}：所有惯性系对于描述力学定律来说都是等价的，力学规律在伽利略变换下形式不变。

    \subsection{狭义相对论基本原理}
    两条公设：狭义相对性原理（物理定律在所有惯性系中形式一致）、光速不变原理（真空中光速对任意惯性系恒为$c$）。

    \subsection{洛伦兹变换}
    \begin{enumerate}
        \item 洛伦兹坐标变换
        \[
        \left\{
        \begin{aligned}
        x' &= \frac{x - vt}{\sqrt{1 - \frac{v^2}{c^2}}}, \\
        y' &= y, \\
        z' &= z, \\
        t' &= \frac{t - \frac{v}{c^2} x}{\sqrt{1 - \frac{v^2}{c^2}}}.
        \end{aligned}
        \right.
        \]
        \item 洛伦兹速度变换
        \[
        \left\{
        \begin{aligned}
        u'_x &= \frac{u_x - v}{1 - \frac{u_x v}{c^2}}, \\
        u'_y &= \frac{u_y \sqrt{1 - \beta^2}}{1 - \frac{u_x v}{c^2}}, \\
        u'_z &= \frac{u_z \sqrt{1 - \beta^2}}{1 - \frac{u_x v}{c^2}}.
        \end{aligned}
        \right.
        \]
        洛伦兹速度逆变换：
        \[
        \left\{
        \begin{aligned}
        u_x &= \frac{u'_x + v}{1 + \frac{u'_x v}{c^2}}, \\
        u_y &= \frac{u'_y \sqrt{1 - \beta^2}}{1 + \frac{u'_x v}{c^2}}, \\
        u_z &= \frac{u'_z \sqrt{1 - \beta^2}}{1 + \frac{u'_x v}{c^2}}.
        \end{aligned}
        \right.
        \]
        其中 $\beta=\dfrac{v}{c}$。
    \end{enumerate}

\section{狭义相对论时空观}
    \subsection{“同时”的相对性}
    不同惯性系下异地事件的同时性不具有绝对性。
    \subsection{时间的相对性}
    运动参考系时间流逝变慢（时间膨胀），固有时最短。
    \subsection{长度的相对性}
    运动方向上物体发生长度收缩，固有长度最长。

\section{相对论动力学基础}
    \subsection{相对论的质速关系}
    运动质量 $m=\dfrac{m_0}{\sqrt{1-v^2/c^2}}$，$m_0$ 为静止质量。
    \subsection{相对论动力学的基本方程}
    $\boldsymbol{F}=\dfrac{\mathrm{d}(m\boldsymbol{v})}{\mathrm{d}t}$，动量定义为 $\boldsymbol{p}=m\boldsymbol{v}$。
    \subsection{相对论动能}
    $E_\mathrm{k}=mc^2-m_0c^2$。
    \subsection{静能、总能和质能关系}
    质能方程 $E=mc^2$，静止物体静能 $E_0=m_0c^2$。
    \subsection{能量和动量的关系}
    相对论能量动量关系式 $E^2 = E_0^2 + (pc)^2$。

% ========== 第二篇 电磁学 ==========
\part{第二篇 电磁学}
% 4.1 第六章 电荷与电场.tex
\chapter{第六章 电荷与电场}

\section{库仑定律与电场强度}
    \subsection{电荷及其性质}
    \subsection{库仑定律}
    \[
    F=k\frac{q_1q_2}{r^3}\boldsymbol{r}
    \]
    比例系数$k$写成
    \[
    k=\frac{1}{4\pi\varepsilon_0}
    \]
    式中，$\varepsilon_0$称为真空介电常量，也叫真空电容率。
    %\[
    %\varepsilon_0\approx 8.85\times10^{-12}\,\mathrm{F/m}
    %\]
    \subsection{电场与电场强度}
    \subsection{场强叠加原理}

\section{电场强度通量与高斯定理}
    \subsection{电场线}
    \subsection{电场强度通量}
    通过某个闭合曲面$S$的$\boldsymbol{E}$通量
    \subsection{真空中高斯定理}
    在真空中静电场内通过任意闭合曲面的$\boldsymbol{E}$通量，等于该闭合曲面所包围的电荷的代数和除以，其数学表达式为
    %\[
    %\]
    若场源电荷连续分布，则式为
    %\[
    %\]
    \subsection{高斯定理的应用}

\section{静电场的环路定理与电势}
    \subsection{静电场的环路定理}
    在静电场中，场强沿闭合路径的线积分等于零，称静电场环路定理。
    \subsection{电势}
    \subsection{电势差}
    \subsection{电势的计算}
    \subsection{等势面 电势梯度}
        \subsubsection{等势面}
        \subsubsection{电势梯度}

\section{静电场中导体}
    \subsection{导体的静电平衡性质}
        \subsubsection{导体的静电平衡状态}
        \subsubsection{静电平衡下导体带电特性}
    \subsection{空腔导体和静电屏蔽}
        \subsubsection{空腔导体}
        \subsubsection{静电屏蔽}
    \subsection{电容和电容器}
        \subsubsection{孤立导体的电容}
        \subsubsection{电容器及其电容}
        \subsubsection{串联和并联}
            \subsubsubsection{串联}
            \subsubsubsection{并联}

\section{静电场中的电介质 电介质中的高斯定理}
    \subsection{电介质的微观结构和极化过程}
    \subsection{电极化强度矢量}
    \subsection{电介质中的场强}
    \subsection{电介质中的高斯定理 电位移\texorpdfstring{$\boldsymbol{D}$}{D}}
    \subsection{有电介质时静电场的计算}

\section{电场的能量}
    \subsection{点电荷系统的电能}
    \subsection{电容器的能量}
    \subsection{电场能量 电场能量密度}
    \subsection{电场能量的计算}
    电场能量为
    %\[
    %d
    %\]
    即可求得整个电场的总能量为
    %\[
    %W_e
    %\]
    显然

\endinput
% 4.2 第七章 电流与磁场.tex
\chapter{第七章 电流与磁场}

\section{恒定电流和恒定电场 电动势}
    \subsection{电流形成的条件}
    定义单位时间内通过导线任一截面的电荷量为电流。

    在导体中形成电流的条件是：（1）导体（1）
    \subsection{恒定电流和恒定电场}
    \subsection{电流和电流密度}
    定义单位时间内通过垂直于电流方向单位面积的电量为导体中某点电流密度矢量的$\boldsymbol{j}$大小，$\boldsymbol{j}$的方向与正电荷在该点
    \subsection{欧姆定律的微分形式 焦耳—楞次定律的微分形式}
    \subsection{电源及电源电动势}
    电源电动势在量值

\section{恒定磁场和磁感应强度}
    \subsection{磁性起源于电荷的运动}
    \subsection{磁场 磁感应强度\texorpdfstring{$\boldsymbol{B}$}{B}}

\section{毕奥‑萨伐尔定律}
    \subsection{毕奥‑萨伐尔定律}
    在真空中
    \subsection{毕奥‑萨伐尔定律应用举例}
    \subsection{匀速运动电荷的磁场}

\section{磁场中的高斯定理}
    \subsection{磁感应线}
    \subsection{磁通量}
    \subsection{恒定磁场中的高斯定理}

\section{真空中恒定磁场的安培环路定理}
    \subsection{恒定磁场的安培环路定理}
    \subsection{安培环路定理的应用}

\section{磁场对运动电荷和载流导线的作用}
    \subsection{洛伦兹力}
    \subsection{带电粒子在磁场中的运动}
    \subsection{应用电场和磁场控制带电粒子运动的实例}
    \subsection{安培力}

\section{磁力的功}
    \subsection{磁力对运动载流导线的功}
    \subsection{磁力矩对转动载流线圈的功}
    磁力矩对转动载流线圈的功等于电流强度与回路包围面积内的磁通量的增量的乘积。

\section{磁介质中的恒定磁场}
    \subsection{磁介质及其磁化}
    \subsection{磁介质中的高斯定理和安培环路定理}
    \subsection{铁磁质}

\endinput
% 4.3 第八章 电磁场与麦克斯韦电磁场方程组.tex
\chapter{第八章 电磁场与麦克斯韦电磁场方程组}

\section{电磁感应的基本定律}
    当穿过闭合导体回路的磁通量发生变化时，闭合导体回路中就会出现电流，称之为电磁感应。

    所以

    导体回路中

    应该

    如果，则
    %\[
    %\]
    在$t_1$到$t_2$这段时间为
    %\[
    %\]
    式表明

    总电动势为
    %\[
    %\]

\section{动生电动势}
    把感应电动势
    \subsection{动生电动势的产生原因}
    \subsection{动生电动势的计算}

\section{感生电动势和感生电场}
    \subsection{感生电动势}
    \subsection{感生电场}
    \subsection{感生电动势的计算}
    \subsection{涡电流}
        \subsubsection{热效应}
        \subsubsection{机械效应}
        \subsubsection{趋肤效应}

\section{自感和互感}
    \subsection{自感}
    \subsection{互感}

\section{磁场的能量}

\section{位移电流和全电流定律}
    \subsection{位移电流}
    \subsection{全电流定律}

\section{麦克斯韦电磁场方程组}

\endinput

% ========== 第三篇 振动与波动 ==========
\part{第三篇 振动与波动}
% 5.1 第十一章 振动学基础.tex
\chapter{第十一章 振动学基础}

\section{简谐运动的描述}
    \subsection{弹簧振子}
    \subsection{简谐运动的表达式}
    \subsection{简谐运动的速度和加速度}
    \subsection{振动的相位}
    \subsection{旋转矢量表示法}

\section{简谐振动的动力学特征}
    \subsection{简谐振动的动力学定义}
    \subsection{简谐振动的实例}
    \subsection{简谐振动的能量}

\section{简谐振动的合成}
    \subsection{两个同方向同频率简谐振动的合成}
    \subsection{两个同方向不同频率简谐振动的合成}
    \subsection{相互垂直的简谐振动的合成}
    质点的轨迹方程为
    %\[
    %\frac{x^2}{A_1^2}+\frac{y^2}{A_2^2}-\frac{2xy}{A_1A_2}\cos\Delta\varphi=\sin^2\Delta\varphi
    %\]

\section{阻尼振动}

\section{受迫振动和共振}

\section{电磁振荡}

\endinput
% 5.2 第十二章 波动学基础.tex
\chapter{第十二章 波动学基础}

\section{机械波的产生和传播}
    \subsection{机械波的形成}
    \subsection{描述波动的物理量}

\section{平面简谐波的波函数}
    \subsection{波函数的建立}
    写成下列形式：
    \[
    y(x,t)=A\cos\left[2\pi\left(\frac{t}{T}-\frac{x}{\lambda}\right)\right]
    \]
    我们，其定义是
    %\[
    %\]
    它可以理解为，波函数式又可写成
    %\[
    %\]
    如果

    以上这些，且，至于，即已知的振动表达式为
    %\[
    %y()
    %\]
    则相应的在$x$轴上传播的简谐波的波函数为
    \[
    y(x,t)=A\cos{\left[\omega\left(t+\frac{x-x_0}{u}\right)+\varphi\right]}
    \]
    \[
    y = A \cos\left[\omega t + \varphi_0 \mp \frac{2\pi}{\lambda} (x - d)\right]
    \]
    \subsection{波函数的物理意义}

\section{波动方程与波速}
    \subsection{物体的弹性形变}
        \subsubsection{线变}
        \subsubsection{切变}
    \subsection{波动方程}
    \subsection{决定波速的因素}

\section{波的能量}
    \subsection{平面简谐纵波传播时介质元的能量}
    仍，当
    %\[
    %\]
    在介质，速度为
    %\[
    %\]
    则介质元的振动动能为
    %\[
    %dE_k=\frac{1}{2}\rho dV\omega^2 A^2\sin^2\left[\omega\left(t-\frac{x}{u}\right)\right]
    %\]

    介质，故弹性势能可写成
    %\[
    %\]
    \subsection{波的能量密度和能流密度}
    介质，可得
    %\[
    %\]
    式说明，因为，所以能量密度的平均值为
    \[
    \overline{w}=\frac{1}{T}\int^T_0 w\mathrm{d}t=\frac{1}{2}\rho A^2 \omega^2
    \]
    式表示波的

    因为，数值$I$为
    \[
    I=\frac{\omega uS}{S}=u=\frac{1}{2}\rho A^2\omega^2 u
    \]
    能流密度

    由于
    \subsection{波的吸收}
    即
    %\[
    %\]
    式，可得
    %\[
    %\]
    式，为
    %\[
    %\]
    \subsection{球面波}

\section{惠更斯原理}
    \subsection{惠更斯原理}
    \subsection{惠更斯原理的应用}

\section{波的叠加原理和波的干涉}
    \subsection{波的叠加原理}
    \subsection{波的干涉}

\section{驻波}
    \subsection{驻波的形成}
    \subsection{驻波的波函数}
    下面，则它们的波函数分别为
    %\[
    %y_1=
    %\]
    %\[
    %y_2=
    %\]
    两波叠加后，合位移为
    \[
    y=y_1+y_2=A\cos\left[2\pi\left(\nu t-\frac{x}{\lambda}\right)\right]+A\cos\left[2\pi\left(\nu t+\frac{x}{\lambda}\right)\right]
    \]
    利用得
    \[
    y=\left(2A\cos{\frac{2\pi}{\lambda}x}\right)\cos{2\pi \nu t}
    \]
    这就是

    从

    \subsection{半波损失}
    波，当波从波疏介质垂直入射波密有半波损失。
    \subsection{弦线振动的简正模式}

\section{多普勒效应}
    现以，$u$表示声波在介质中的传播速度

    设波源频率为，下面分三种情况讨论：

    (1) 波源$S$静止在介质中，观测者相对于介质以速度$v_0$向着波源移动，观测者向着波源运动时，感到波是以速度为，即观测者测得的频率为
    %\[
    %\]

    (2) 设观测者静止在介质中，看出
    %\[
    %\]
    所以观测者每秒内接收到的完整波的个数，即频率$\nu'$为
    \[
    \nu'=\frac{u}{\lambda'}=\frac{u}{(u-u_s)T}=\nu\cdot\frac{u}{u-u_s}
    \]
    当

    (3) 设波源，为
    \[
    \nu'=\nu\left(\frac{u+u_0}{u-u_s}\right)
    \]

\section{声波}

\section{电磁波}

\endinput

\end{document}